\documentclass[../../main.tex]{subfiles}
\begin{document}
\section{Family 4: moment maps, Monge--Amp\`ere, and Stein kernels}
\label{sec:family-moment-map}

\paragraph{Object followed.}
The Hessian $H=\Hess\varphi$ of the moment-map potential of $\mu$, in the coordinates in which
$\mu$ is the pushforward of its own moment measure.

\paragraph{What it buys.}
This is where the July 2026 progress happened, and it is currently the cleanest new proof
component in the subject. In moment-map coordinates the statement ``$\mu$ is isotropic'' becomes
``$\E_\nu H=I$'', and the Monge--Amp\`ere equation, differentiated twice, produces two
manifestly positive semidefinite source terms. Exploiting that positivity yields a sharp bound on
$\E\Tr(BHBH)$ for every \emph{constant} symmetric matrix $B$ --- from which the sharp thin-shell
constant, the sharp third-moment tensor, and the all-quadratic Poincar\'e inequality all follow.

This section presents the mechanism. The repository's formal statements of the two consequences it
actually consumes are Theorem~\ref{thm:letwin-qcts} (Section~\ref{sec:qcts}) and
Proposition~\ref{prop:letwin-kappa} (Section~\ref{sec:covariance-tech}); they are not restated
here.

\subsection{Moment-map coordinates}
\label{subsec:mm-coordinates}

For a centered, full-dimensional $\mu=e^{-V}\dd x$, the moment-measure theorem supplies an
essentially unique convex $\varphi$ with
\begin{equation}
  (\nabla\varphi)_\#\nu=\mu,
  \qquad
  \dd\nu(y)=e^{-\varphi(y)}\dd y
  \label{eq:moment-measure}
\end{equation}
\cite{CorderoErausquinKlartag2015MomentMeasures,Klartag2013MomentMeasures}. Set
\[
  H(y)=\Hess\varphi(y).
\]
Change of variables in \eqref{eq:moment-measure} is the Monge--Amp\`ere equation
\begin{equation}
  e^{-V(\nabla\varphi(y))}\det H(y)=e^{-\varphi(y)},
  \qquad\text{that is}\qquad
  \log\det H=-\varphi+V(\nabla\varphi).
  \label{eq:MA}
\end{equation}

If $Y\sim\nu$ and $X=\nabla\varphi(Y)$, then $X\sim\mu$, and integration by parts under $\nu$ gives
\begin{equation}
  \E_\nu H(Y)=\E_\nu\bigl[\nabla\varphi(Y)\otimes\nabla\varphi(Y)\bigr]=\Cov\mu=I
  \label{eq:mean-hessian-identity}
\end{equation}
when $\mu$ is isotropic. So \emph{``average Hessian equals identity'' replaces isotropy} in
moment-map coordinates. This is the structural payoff: isotropy, which
Section~\ref{sec:family-needles} showed is not inherited by needles, becomes a single linear
identity for the object being estimated.

\begin{theorem}[Compact-target regular moment map]
\label{thm:regular-moment-map-compact-target}\klsstatus{thm:regular-moment-map-compact-target}
Let $P\subset\R^n$ be a convex body and let
$\mu(dx)=g(x)\mathbf 1_{\operatorname{int}P}(x)\,dx$ be a centered probability measure, where
$g\in C^\infty(\R^n)$ is positive. Then the canonical moment potential $\varphi$ is smooth and
strictly convex on $\R^n$, and
\[
  \nabla\varphi:\R^n\longrightarrow\operatorname{int}P
\]
is a diffeomorphism. Moreover,
\[
  \tau_\mu(x)=D^2\varphi\bigl((\nabla\varphi)^{-1}(x)\bigr)
\]
is a smooth positive symmetric Stein kernel: $\Div_\mu\tau_\mu=-x$ distributionally, with weak
zero boundary flux, and $\E_\mu\tau_\mu=\Cov(\mu)$.
\end{theorem}

The moment-potential regularity and global gradient-image statement are
Theorem~1.1 of \cite{BermanBerndtsson2013RealMA}; the target-coordinate Stein identity and weak
zero-flux formulation are Theorem~2.3 of \cite{Fathi2019SteinMomentMaps}. This compact-target
class is the published regularity input used by the approximation closure in
Question~\ref{q:cmh-approximation}.

\subsection{Differentiating Monge--Amp\`ere}
\label{subsec:mm-differentiated}

Introduce the elliptic operator associated with the Hessian metric,
\begin{equation}
  \calL f=\varphi^{ij}\partial_{ij}f-(V_i\circ\nabla\varphi)\,\partial_if,
  \qquad (\varphi^{ij})=H^{-1},
  \label{eq:calL-def}
\end{equation}
which is symmetric in $L^2(\nu)$:
\begin{equation}
  \E_\nu[(\calL f)g]=-\E_\nu\bigl[\varphi^{ij}(\partial_if)(\partial_jg)\bigr].
  \label{eq:calL-symmetric}
\end{equation}
Differentiating \eqref{eq:MA} twice gives the key identity
\begin{equation}
  \calL H_{ij}+H_{ij}
  =\bigl(H\,\Hess V(\nabla\varphi)\,H\bigr)_{ij}
  +\varphi^{ac}\varphi^{bd}\varphi_{abi}\varphi_{cdj} .
  \label{eq:differentiated-MA}
\end{equation}
Both terms on the right are positive semidefinite: the first by convexity of $V$, the second
because it is a Gram matrix of normalized third derivatives. \emph{This positivity is the PDE
source of the new estimate}; everything below is a way of harvesting it.

\subsection{The fixed-matrix Hessian bound}
\label{subsec:mm-fixed-matrix}

Let $B\succeq0$ and define
\[
  S_B=\Tr(BHBH),
  \qquad
  D_B=\varphi^{ab}\Tr\bigl(B(\partial_aH)B(\partial_bH)\bigr).
\]
Applying $\calL$ to $S_B$, integrating against $\nu$, and using the two positive sources of
\eqref{eq:differentiated-MA} yields
\begin{equation}
  \E S_B\ge2\,\E D_B .
  \label{eq:SB-DB}
\end{equation}
The nontrivial point is a third-derivative comparison. At a fixed point, normalize $H=I$ and
diagonalize $B=\diag(b_i)$; writing $T_{ijk}=\varphi_{ijk}$, the difference between the two sides
reduces to
\begin{equation}
  \tfrac12\sum_{i,j,k}(b_i-b_j)^2T_{ijk}^2\ \ge\ 0 .
  \label{eq:third-derivative-comparison}
\end{equation}

Next apply Brascamp--Lieb entrywise to $F=B^{1/2}HB^{1/2}$. Since $\E F=B$ by
\eqref{eq:mean-hessian-identity},
\[
  \E\norm{F-B}_{\HS}^2=\E S_B-\Tr(B^2),
\]
and the Brascamp--Lieb energy of $F$ is exactly $D_B$, so
\[
  \E S_B-\Tr(B^2)\ \le\ \E D_B\ \le\ \tfrac12\E S_B ,
\]
the last step by \eqref{eq:SB-DB}. Rearranging gives the main new theorem inside Letwin's proof.

\begin{theorem}[Letwin matrix moment-map estimate]
\label{thm:letwin-moment-map}\klsstatus{thm:letwin-moment-map}
Under the regular moment-map assumptions used in \cite{Letwin2026QuadraticKLS}, for every
constant symmetric matrix $B$,
\begin{equation}
  \E_\nu\Tr\bigl(BH(Y)BH(Y)\bigr)\le2\Tr(B^2).
  \label{eq:letwin-matrix}
\end{equation}
\end{theorem}

For indefinite $B$ the statement still holds: diagonalizing gives the pointwise bound
$\Tr(BHBH)\le\Tr(\abs BH\abs BH)$, and \eqref{eq:letwin-matrix} applies to $\abs B\succeq0$.

The $B=I$ case of \eqref{eq:letwin-matrix} is the moment-Hessian estimate shared with the
contemporaneous Chen--Klartag preprint.

\begin{theorem}[Chen--Klartag moment-Hessian estimate]
\label{thm:chen-klartag-moment-hessian}\klsstatus{thm:chen-klartag-moment-hessian}
Under the regularity assumptions of \cite{ChenKlartag2026SharpThinShell},
\[
  \E_\nu\norm H_{\HS}^2\le2n.
\]
The general log-concave conclusions below follow by the approximation argument in that preprint.
\end{theorem}

\begin{theorem}[Chen--Klartag sharp thin shell]
\label{thm:chen-klartag-thin-shell}\klsstatus{thm:chen-klartag-thin-shell}
For every isotropic log-concave $X\in\R^n$,
\[
  \Var(|X|^2)\le8n.
\]
The constant is attained by products of standard centered one-sided exponential variables.
\end{theorem}

\begin{theorem}[Chen--Klartag sharp third tensor]
\label{thm:chen-klartag-third-moment}\klsstatus{thm:chen-klartag-third-moment}
If $X\in\R^n$ is isotropic and log-concave and $T_3(X)=(\E X_iX_jX_k)_{i,j,k}$, then
\[
  \norm{T_3(X)}_{\HS}^2\le4n.
\]
The preprint also proves sharper convex-body estimates with equality for the regular simplex.
\end{theorem}

\begin{remark}[What each preprint contributes]
\label{rem:mm-two-preprints}
Theorems~\ref{thm:letwin-moment-map}--\ref{thm:chen-klartag-third-moment} are imported from
contemporaneous version-1 preprints
\cite{ChenKlartag2026SharpThinShell,Letwin2026QuadraticKLS}; they are not promoted here as
independently certified repository proofs. Letwin's $B=I$ case contains the shared moment-Hessian
estimate, while Chen--Klartag's sharp third-tensor, cone, simplex, and equality analysis is
additional. Conversely, Letwin controls every constant symmetric homogeneous quadratic form.
Neither subsumes the other. Only Letwin's preprint supplies the claimed improvement of the general
KLS bound.
\end{remark}

\subsection{The Stein kernel and the \texorpdfstring{$H^{-1}$}{H-1} inequality}
\label{subsec:mm-stein}

Fathi observed that the moment-map Hessian, read in target coordinates,
\begin{equation}
  \tau_\mu(x)=H\bigl((\nabla\varphi)^{-1}(x)\bigr),
  \label{eq:stein-kernel-def}
\end{equation}
is a positive symmetric Stein kernel for $\mu$ \cite{Fathi2019SteinMomentMaps}:
\begin{equation}
  \E[X_if(X)]=\E\sum_j\tau_{ij}(X)\,\partial_jf(X).
  \label{eq:stein-identity}
\end{equation}
Combining \eqref{eq:stein-identity} with the $H^{-1}$ calculus of Section~\ref{subsec:hminus1}
gives, for a linear function $\ell_v(x)=\inner xv$,
\begin{equation}
  \norm{\ell_v}_{H^{-1}(\mu)}^2\le\E\abs{\tau_\mu(X)v}^2 .
  \label{eq:stein-hminus1}
\end{equation}

By Remark~\ref{rem:quadratics-special}, quadratic test functions have linear derivatives, so
\eqref{eq:stein-hminus1} can be fed into the Barthe--Klartag inequality
\eqref{eq:barthe-klartag}. This is the junction at which the two families meet.

\subsection{The noncommutativity trick}
\label{subsec:mm-noncommutativity}

There is one genuine obstacle between \eqref{eq:letwin-matrix} and a quadratic Poincar\'e
inequality. For $q_M(x)=\inner{Mx}x$, applying \eqref{eq:stein-hminus1} directly produces
\[
  \E\norm{\tau_\mu(X)M}_{\HS}^2 ,
\]
which is \emph{not} the expression controlled by \eqref{eq:letwin-matrix}, because $M$ and
$\tau_\mu(X)$ need not commute.

The resolution is a change of variables. Write $M=\operatorname{sgn}(M)\abs M$ and set
$Z=\abs M^{1/2}X$, with $\eta$ the law of $Z$. Then
\[
  q_M(X)=\inner{\operatorname{sgn}(M)Z}Z ,
\]
and the Stein kernel transforms by congruence,
\begin{equation}
  \tau_\eta(Z)=\abs M^{1/2}\,\tau_\mu(X)\,\abs M^{1/2},
  \label{eq:stein-congruence}
\end{equation}
so that
\begin{equation}
  \E\norm{\tau_\eta(Z)}_{\HS}^2=\E\Tr\bigl(\abs MH\abs MH\bigr)\le2\Tr(M^2)
  \label{eq:conjugated-bound}
\end{equation}
by Theorem~\ref{thm:letwin-moment-map} applied with $B=\abs M$. The conjugation has moved $M$
inside the trace in exactly the pattern \eqref{eq:letwin-matrix} controls.

Applying the $H^{-1}$ inequality to $F(z)=\inner{\operatorname{sgn}(M)z}z-\Tr M$ and using that
$\operatorname{sgn}(M)$ is orthogonal gives
\[
  \Var\inner{MX}X\le4\,\E\norm{\tau_\eta(Z)}_{\HS}^2\le8\Tr(M^2),
\]
while isotropy gives $\E\abs{\nabla\inner{MX}X}^2=4\Tr(M^2)$. Together these are the quadratic
Poincar\'e inequality recorded as Theorem~\ref{thm:letwin-qcts}, with the constant $2$ that is
sharp for products of centered exponentials at $M=I$.

The reduction from there to $\kappa_n\le2\sqrt2$ is short and purely algebraic; it is carried out
as Proposition~\ref{prop:letwin-kappa}. Substituting into the bridge \eqref{eq:kls-bridge} gives
the current preprint record.

The constant $2\sqrt2$ is not sharp, and Route~C says exactly what would sharpen it. On the
regular moment-map class, Corollary~\ref{cor:gate-zero-third-moment} turns any bound
$\E[\tau^2]\preceq c\,\Id$ into the directional third-moment bound
$\norm{T_3(a)}_{\HS}\le2\sqrt{c-1}$, so sharp gate zero at $c=2$
(Conjecture~\ref{conj:gate-zero-sharp}) would give the sharp $\kappa_n\le2$, attained by
products of centered exponentials. The implication runs one way only: a sharp third-moment
bound does not return sharp gate zero, because the high-mode remainder of
Lemma~\ref{lem:linear-sector-third-moment} is not zero off the cone axis.

\subsection{Audit and epistemic status}
\label{subsec:mm-audit}

Letwin's paper is an arXiv version-1 preprint of 27 July 2026, not a peer-reviewed result, and it
carries an explicit AI-use disclosure. The identities above are internally coherent, and the
algebra from the quadratic theorem to $\kappa_n\le2\sqrt2$ is transparent. The delicate points for
independent checking are concentrated elsewhere:

\begin{enumerate}[label=(\alph*),leftmargin=2.2em]
  \item the coordinate invariance of the third-derivative comparison
  \eqref{eq:third-derivative-comparison};
  \item noncompact integration by parts, specifically the vanishing $\E\calL S_B=0$ used to obtain
  \eqref{eq:SB-DB};
  \item approximation and moment-map regularity, that is, the passage from the regular case to
  arbitrary log-concave laws;
  \item the extraction of $\CP\lesssim\kappa_n\sqrt{\log n}$, which the preprint obtains by
  \emph{tracing} Klartag's estimates \cite{Klartag2023Logarithmic} rather than by quoting a
  theorem stated in that form.
\end{enumerate}

The preprint's appendix addresses each of (a)--(c). Point (d) is a concatenation of Klartag's
variance-transfer result with a precise short-time covariance estimate, and is the step this
repository re-derives for its own use in Section~\ref{sec:covariance-tech}.

Accordingly, throughout Parts~III this result is described as the best \emph{current} bound and
Klartag's as the best \emph{published} bound. Every repository node depending on it carries
\texttt{import\_class: preprint-unreviewed}, which by the ledger contract cannot underlie a node
marked \texttt{proved}. The same status distinction is drawn elsewhere in the literature; see
\cite{Zhang2026HitAndRun}, whose bibliographic metadata in this repository is itself flagged as
unverified.

\paragraph{The precise missing estimate.}
Control of $\E\inner{\tau_\mu(X)\nabla f(X)}{\nabla f(X)}$ for an arbitrary $f$, in place of
$\E\Tr(BHBH)$ for a constant matrix $B$.

\paragraph{Why it stalls.}
Brascamp--Lieb in moment-map coordinates already gives
\begin{equation}
  \Var_\mu f\le\E_\mu\inner{\tau_\mu(X)\nabla f(X)}{\nabla f(X)},
  \label{eq:mm-brascamp-lieb}
\end{equation}
so KLS would follow from $\E\inner{\tau_\mu\nabla f}{\nabla f}\lesssim\E\abs{\nabla f}^2$. The
identity $\E\tau_\mu=I$ of \eqref{eq:mean-hessian-identity} is \emph{not} sufficient for this,
because $\tau_\mu(X)$ can correlate with $\nabla f(X)$. Letwin controls a deterministic $B$; the
missing theorem must handle an $X$-dependent matrix or direction field.

\paragraph{How this family feeds the live routes.}
It supplies two of them. Route~C (Section~\ref{sec:moment-map-cmh}) attacks the gap above
head-on, replacing constant multipliers by test-dependent Haar fields, and its endpoint
$\CMH$ is defined in these coordinates. Route~S consumes the family's quadratic control as the
estimate it feeds its whitened posterior tensor to. The fixed-matrix bound
(Theorem~\ref{thm:letwin-moment-map}) is therefore the single literature input both routes are
trying to make adaptive.
\end{document}
